\begin{resumo}[Abstract]
\begin{otherlanguage*}{english}
 
Biodiesel production by alkaline transesterification faces difficulties related to the low miscibility between the triglyceride and the alcohol used in the reaction, which tends to compromise process efficiency and the subsequent phase separation step. In this context, protic ionic liquids (PILs) have been investigated as cosolvents capable of modulating the liquid-liquid equilibrium (LLE) behaviour of these systems, although the mechanisms behind such modulation remain little explored in the literature, particularly for ternary systems involving triglycerides, for which experimental LLE data are unavailable. This work aimed to thermodynamically model the LLE of twelve ternary systems at 298.15~K, combining three triglycerides of increasing unsaturation (triolein, trilinolein and trilinolenin), two alcohols (ethanol and methanol) and two PILs derived from the hexanoate anion, 2-hydroxyethylammonium hexanoate ([2-HEA][Hx]) and diethylenetriammonium hexanoate ([DETA][Hx]), assessing through two predictive approaches of distinct nature whether these compounds increase or reduce the miscibility of the oil-alcohol pair: UNIFAC-LL, a group-contribution model, and COSMO-SAC, an \textit{a priori} model based on surface charge density profiles obtained from quantum chemical calculations. The interaction parameters missing from the UNIFAC-LL matrix for the groups characteristic of the PILs were supplied by a machine learning model, and the calculation engine common to both models was externally verified against reference systems from the literature. In the eighteen systems for which a binodal curve was obtained, without exception, the width of the immiscibility gap increases with the PIL mole fraction, with slopes between $+0.87$ and $+2.95$, so that the cosolvency hypothesis is not confirmed under the modelled conditions: the PIL drives the coexisting phases apart rather than bringing them together, behaving as a separating agent. In both models the widening arises from the displacement of the oil-rich branch, whereas the composition of the alcohol-rich phase remains virtually unchanged, which suggests that the yield gains reported experimentally, obtained at different temperature and concentration, are not explained by homogeneous phase solubilisation. The models agree on the systems containing [2-HEA][Hx], for which both indicate a two-phase region, and disagree on those containing [DETA][Hx], where COSMO-SAC predicts phase separation in all six cases while UNIFAC-LL does not allow the curve to be constructed in any of them, a disagreement that the available data cannot settle in the absence of experimental measurements. The physicochemical interpretation rests on the anchoring of the PIL to the alcohol-rich phase through hydrogen bonding, the resulting competition for the alcohol hydroxyl group and the steric hindrance of the chains esterified to the glycerol backbone.
 
\vspace{\onelineskip}
\noindent
\textbf{Keywords}: Protic ionic liquids. Biodiesel. Liquid-liquid equilibrium. UNIFAC-LL. COSMO-SAC.
 
\end{otherlanguage*}
\end{resumo}